% !TeX root = ../main.tex
\section{模型假设}

\subsection{问题一的抽样检测假设}

表~\ref{tab:q1-assumptions}列出问题一采用的假设、建模作用及适用范围。每项假设均直接影响抽样分布或结论解释；样本量口径、检验功效等建模边界见 5.1 节开头说明。

\begin{table}[H]
  \centering
  \caption{问题一模型假设及其作用}
  \label{tab:q1-assumptions}
  \begin{tabularx}{\textwidth}{c>{\raggedright\arraybackslash}p{0.42\textwidth}>{\raggedright\arraybackslash}X}
    \toprule
    编号 & 假设 & 建模作用与适用范围 \\
    \midrule
    A1 & 每件零配件只能判为合格或次品，且检测结果准确。 & 将观测建模为伯努利变量；若存在误检或漏检，需另行引入检测灵敏度与特异度。 \\
    A2 & 总体量未知或抽样比例较小时，给定真实次品率 $p$ 后各次检测相互独立且同分布。 & 允许使用二项分布；大比例无放回抽样不适用该假设。 \\
    A3 & 有限总体分支中，批次总量 $N$ 已知，样品从同一批次均匀且无放回抽取。 & 允许使用超几何分布进行有限总体修正。 \\
    \bottomrule
  \end{tabularx}
\end{table}

\subsection{问题二的生产决策假设}

表~\ref{tab:q2-assumptions}列出问题二采用的假设及建模作用。

\begin{table}[H]
  \centering
  \caption{问题二模型假设及其作用}
  \label{tab:q2-assumptions}
  \begin{tabularx}{\textwidth}{c>{\raggedright\arraybackslash}p{0.44\textwidth}>{\raggedright\arraybackslash}X}
    \toprule
    编号 & 假设 & 建模作用 \\
    \midrule
    B1 & 给定质量参数后，各次采购与装配结果相互独立，次品率与成本在一个交付周期内保持不变。 & 使状态转移与一步成本保持齐次，MDP 可解。 \\
    B2 & 检测结果准确，检测出的不合格零配件被丢弃。 & 检测动作将状态确定地更新为已知合格或缺失。 \\
    B3 & 拆解不损伤零配件，拆解前已确认合格的零配件仍保持已知合格。 & 保留拆解前的质量信息，避免重复检测。 \\
    B4 & 未检测零配件的真实质量不可直接观察，仅以联合信念表示。 & 需要联合后验而不是两个独立边际概率。 \\
    B5 & 采购、检测、装配、市场调换与拆解成本均按题目口径计入，交付一件合格成品后过程终止。 & 定义一步成本与终止条件，期望利润等于售价减初始期望成本。 \\
    \bottomrule
  \end{tabularx}
\end{table}

\subsection{问题三的生产决策假设}

问题三沿用问题二假设，并针对多工序结构补充表~\ref{tab:q3-assumptions}中的约定。

\begin{table}[H]
  \centering
  \caption{问题三模型假设及其作用}
  \label{tab:q3-assumptions}
  \begin{tabularx}{\textwidth}{c>{\raggedright\arraybackslash}p{0.44\textwidth}>{\raggedright\arraybackslash}X}
    \toprule
    编号 & 假设 & 建模作用 \\
    \midrule
    C1 & 零配件没有主动报废或拆解决策，检测出的不合格零配件被丢弃后重复购检。 & 规定零配件层面的唯一处置方式。 \\
    C2 & 拆解决策只作用于不合格半成品或不合格成品，且拆解不损坏下级零部件，回流后保留原真实质量状态、不重新按次品率抽样。 & 使拆解回流成为确定性的状态恢复。 \\
    C3 & 任一下级输入不合格则上级必不合格；仅当全部下级输入合格时，上级才以给定次品率产生装配缺陷。 & 决定上级产品合格概率。 \\
    C4 & 固定策略下各状态的一步成本与转移在给定参数内保持不变。 & 使期望成本满足线性方程组，可精确求解。 \\
    \bottomrule
  \end{tabularx}
\end{table}

\subsection{问题四的参数不确定性假设}

问题四沿用问题二、问题三的生产规则，并补充表~\ref{tab:q4-assumptions}中的参数不确定性假设。其中抽样样本量 $n$ 为证据强度参数，$n=40$ 仅为能够兼容题面各名义次品率的代表性小样本情景，企业提供真实抽样方案后可直接替换。
\begin{table}[H]
  \centering
  \caption{问题四模型假设及其作用}
  \label{tab:q4-assumptions}
  \begin{tabularx}{\textwidth}{c>{\raggedright\arraybackslash}p{0.44\textwidth}>{\raggedright\arraybackslash}X}
    \toprule
    编号 & 假设 & 建模作用 \\
    \midrule
    D1 & 表~1、表~2 的次品率来自固定样本量为 $n$ 的抽样检测，样本次品数为 $k=n\hat p$（取整）；$n$ 为证据强度参数，$n=40$ 为代表性小样本情景。 & 把名义次品率转化为可更新的后验信息；真实 $n$ 确定后需重算。 \\
    D2 & 各环节次品率采用无信息均匀先验 $\operatorname{Beta}(1,1)$，后验为 $\operatorname{Beta}(k+1,n-k+1)$。 & 使后验具有闭式形式，情景抽样无需数值积分。 \\
    D3 & 各环节真实次品率在给定抽样信息下按独立后验建模，同一实验设置下所有策略共享同一组后验情景。 & 保证策略间比较口径一致，避免抽样噪声造成的不公平对比。 \\
    D4 & 抽样不确定性只改变次品率输入，生产流程、成本结构与决策规则沿用问题二、问题三；问题二在七变量固定策略类中重决策，问题三在 16 位策略空间中重决策。 & 使本问可按 5.2 节与 5.3 节相同的质量传递与成本规则做解析期望计算，结果可与点估计口径逐项对照。 \\
    \bottomrule
  \end{tabularx}
\end{table}
